\documentclass[11pt]{article}
\usepackage{fontspec,amsmath,amssymb,amsthm,hyperref}
\usepackage[margin=1in]{geometry}
\setlength{\emergencystretch}{3em}
\title{Transfer Gibbs observations to the actual stationary process}
\author{LangevinLean mathematical exposition}
\date{}
\begin{document}\maketitle
\noindent Source module: \texttt{AverageLMC/LaplacianStationaryTest.lean}. Compile this proof document with XeLaTeX or LuaLaTeX. The accompanying HTML pairs the mathematical argument with the exact Lean statements and records the scope of the formalization.


Stationarity is an equality of actual marginal laws. This module transfers genuine L2 moments and approximation errors from the Gibbs measure to the constructed process, and converts finite coordinate Dynkin identities into vector inner products.

\[\mathcal L(Y)=\pi\quad\Longrightarrow\quad\|F(Y)-G(Y)\|_{\mathbb L_2}^2=\int\|F-G\|^2\,d\pi,\quad \langle\Psi,F(Y)\rangle_{\mathbb L_2}=\sum_j\mathbb E[\Psi_jF_j(Y)].\]
\section{Represent every square-integrable stationary observation}


Rewrite the observation's Gibbs moment with the actual marginal map and pull it back to the original probability space. The resulting MemLp proof defines a genuine quotient-valued observation. Its almost-everywhere representative is exactly the original vector function evaluated at the actual process state.



\noindent Lean declarations: \texttt{stationary\_observation\_memLp}, \texttt{stationaryObservation}, \texttt{stationaryObservation\_coe}.

\section{Preserve the exact squared approximation error}


The L2 inner-product formula identifies the squared norm with the integral of the squared representative norm. The proved representative identity and the stationary law map turn this into the actual Gibbs moment. Applying the same argument to differences preserves the complete squared approximation error, with no change of constants or dependence on time.



\noindent Lean declarations: \texttt{stationaryObservation\_norm\_sq}, \texttt{stationaryObservation\_sub\_norm\_sq}.

\section{Sum scalar weighted identities into vector duality}


Each coordinate product is integrable by domination with the product of the two vector norms and the L2 Cauchy inequality. Expand the Euclidean inner product as the finite coordinate sum, then interchange this finite sum and the genuine integral. The result is exactly the sum of scalar weighted expectations used by the compact-test Dynkin theorem; no separate coordinatewise bound introduces a dimension factor.



\noindent Lean declarations: \texttt{stationaryObservation\_inner\_eq\_sum}.

\end{document}
